\section{Comparison with ReCom-Based Methods}
\label{sec:comparison}

\subsection{The Four-Method Landscape}

The redistricting toolbox now contains four distinct search and sampling
methods, each occupying a different cell in the (cost, exploration radius,
use case) design space.
Table~\ref{tab:comparison} summarises the key dimensions.

\begin{table}[ht]
\centering
\caption{Comparison of redistricting search and sampling methods.
  Acceptance rate and EC range are empirical estimates from our experiments;
  all values are approximate and state-dependent.
  $n$ = number of tracts, $k$ = number of districts, $T$ = convergence
  threshold (default 600 for \cs{}), $F$ = flip steps (default 10{,}000).}
\label{tab:comparison}
\renewcommand{\arraystretch}{1.35}
\begin{tabular}{@{}lllllll@{}}
\toprule
Method & Steps & Step cost & Accept & EC range & Partisan stab. & Use case \\
\midrule
\flip{}
  & $F = 10\text{K}$
  & $O(\deg)$
  & $\approx\!55\%$
  & Narrow
  & High ($\geq\!0.90$)
  & Sensitivity cert. \\
\sburst{}~\citep{g6paper}
  & $1\text{K} \times B$
  & $O(n/k \log n/k)$
  & $\approx\!15\%$
  & Moderate
  & Moderate
  & Optimisation \\
\be{}~\citep{h1paper}
  & $100 \times \text{nodes}$
  & $O(n/k \log n/k)$
  & $\approx\!20\%$
  & Moderate
  & Moderate
  & Structure \\
\cs{}~\citep{b11paper}
  & $T = 600$
  & $O(n \log n / k)$
  & N/A (METIS)
  & Wide
  & Low ($\leq\!0.70$)
  & Global search \\
\bottomrule
\end{tabular}
\end{table}

\subsection{Flip vs.\ Short-Burst}

\sburst{}~\citep{cannon2022,g6paper} runs $B$ short \recom{} chains of length
$L$ (burst length), selecting the best plan found in each burst.
The key difference from \flip{} is the step proposal: \recom{} merges two
districts and resamples a spanning tree, potentially moving $O(n/k)$ tracts
in one step, while \flip{} moves exactly one tract.

The consequences for sensitivity analysis are:
\begin{itemize}
  \item \sburst{} escapes local optima faster---it can tunnel through the
        topological obstacles that trap \flip{} (Theorem~\ref{thm:reachability}).
  \item \sburst{} has a lower acceptance rate ($\approx\!15\%$ vs. $55\%$ for
        \flip{}) because spanning-tree resampling rarely finds a balanced cut
        on the first try.
  \item \sburst{} explores a wider EC range than \flip{} in the same wall-clock
        time, making it better for plan optimisation (finding the minimum-EC
        plan in a neighbourhood) rather than local sensitivity (characterising
        the neighbourhood of a fixed plan).
\end{itemize}

For sensitivity certification, \flip{} is strictly preferable to \sburst{} because
it directly characterises the single-tract neighbourhood of the submitted plan.
\sburst{} is preferable for optimisation tasks where the goal is to find a better
plan, not to certify the submitted plan.

\subsection{Flip vs.\ BisectionEnsemble}

\be{}~\citep{h1paper} runs short \recom{} chains at each node of the
bisection tree generated by \cs{}.
At each node, it samples alternative sub-district configurations that are
compatible with the fixed higher-level splits.
This makes \be{} a \emph{conditional} ensemble: it samples plans that share
the same hierarchical structure as the submitted plan but vary in sub-district
assignment.

The key difference from \flip{} is the level of perturbation:
\be{} can change $O(n/k)$ tracts in a single node's assignment, while
\flip{} changes exactly one tract globally.
\be{} is therefore better for characterising the set of plans with the same
hierarchical structure (useful for demonstrating that the chosen structure is
not uniquely favourable), while \flip{} is better for certifying robustness
to minimal perturbations.

\subsection{Flip vs.\ ConvergenceSweep}

\cs{}~\citep{b11paper} runs $T$ independent METIS seeds and selects the
best plan by edge-cut.
The plans it generates are independent of each other and of the submitted plan;
there is no notion of ``neighbourhood'' in the \cs{} output.

\flip{}, by contrast, generates plans that are within Hamming distance
$\lfloor F \cdot p_{\mathrm{acc}} \rfloor$ of the submitted plan in the
worst case, though the actual Hamming distance of any visited plan is much
smaller because the chain wanders randomly.

The two methods serve complementary purposes:
\cs{} answers ``is there a better plan out there?''
\flip{} answers ``is the submitted plan stable to local perturbations?''
A complete submission uses both: \cs{} to generate the plan, and \flip{}
to certify its local stability before filing.

\subsection{Computational Cost in Practice}

For North Carolina ($n = 2{,}200$, $k = 14$), typical wall-clock times
on a single core are:
\begin{itemize}
  \item \flip{} ($F = 10{,}000$): approximately 1--3 seconds.
  \item \sburst{} ($B = 50$ bursts, $L = 20$ steps): approximately 15--30 seconds.
  \item \be{} (100 chains per node, 28 nodes): approximately 5--15 minutes.
  \item \cs{} ($T = 600$ seeds): approximately 3--8 minutes.
\end{itemize}
\flip{} is the cheapest method by a large margin.
The full $p$-sweep (five \flip{} runs at five percentiles) takes under 15
seconds and provides the legislative disclosure required by the audit protocol
in Section~\ref{sec:audit}.
